%@ pages 322-324
%@ status translated
% Front matter: title page, attribution, translator's preface, contents.

\begin{titlepage}
\centering
\vspace*{1.5cm}

{\large\bfseries ABELIAN CATEGORIES\par}
\vspace{0.6cm}
{\large\itshape Des catégories abéliennes\par}

\vspace{1.8cm}
{\large PIERRE GABRIEL\par}

\vspace{2.4cm}
{\large Bulletin de la Société Mathématique de France\par}
{\large Tome 90 (1962), p. 323--448\par}

\vfill

\begin{minipage}{0.82\textwidth}
\small
\noindent\copyright\ Bulletin de la Société Mathématique de France, 1962.
All rights reserved.

\smallskip
\noindent This is an unofficial English translation of the article cited above,
prepared from the Numdam digitisation
(\url{https://www.numdam.org/item/10.24033/bsmf.1583/}). It is offered for
scholarly use; in case of any discrepancy the original French text of
Gabriel (1962) is authoritative.

\smallskip
\noindent The article was also published as the author's doctoral thesis
(Thèse Sc. math., Paris, 1961).
\end{minipage}

\vspace{1cm}
\end{titlepage}

\chapter*{Preface to the translation}
\addcontentsline{toc}{chapter}{Preface to the translation}

This translation is complete in intent: its aim is to reproduce Gabriel's
argument completely, in English, with its definitions, hypotheses and
conclusions as they stand. Two conventions govern the rendering.

\medskip
\noindent\textbf{Page numbers.}
The body of this document is numbered so that page~323 of the translation
corresponds to page~323 of the original. English prose does not occupy
exactly the same space as French, so the correspondence drifts slowly as the
text proceeds; the section, theorem and page structure is otherwise identical
to the original, and the original page number is recorded in the margin of
each source file in the repository.

\medskip
\noindent\textbf{Terminology.}
Gabriel writes in a vocabulary that was still being invented in 1962, and
which Grothendieck altered substantially within a few years --- for instance
Gabriel writes \emph{foncteur} where the same author would soon write
\emph{foncteur} in a narrower sense, and \emph{catégorie avec générateurs et
limites inductives exactes} for what is now called a Grothendieck category.
Table~\ref{tab:terminology} records, term by term, what has been done. Where
standard present-day English exists it has been used, so that the results can
be found in the modern literature under the name given here; where the
original coinage is itself part of the history --- the concept of a universe,
Serre's language modulo $\mathcal{C}$, the theory of \emph{fondeurs exacts à
gauche} --- the French is given at first occurrence and the modern term is
used afterwards.

\medskip
\noindent\textbf{Definitions.}
Gabriel's definitions are kept in their original form: a list of data followed
by a list of axioms, rather than compressed into a sentence. This makes the
translation checkable against the original clause by clause.

\medskip
\noindent\textbf{Notation.}
Gabriel uses bold upright letters for categories ($\Cat$, $\Ens$, $\Ab$) and
Fraktur for universes; both are reproduced here, so that $\mathbf{C}$ in this
translation is $\mathbf{C}$ on page 330 of the original.

\tableofcontents

% Terminology and editorial conventions
% ========================================
%
% Gabriel's 1962 French coinages do not all survive into modern English, and
% several of them were superseded by Grothendieck's own usage within a few
% years.  The policy for this translation is:
%
%   1. Where standard English mathematical usage exists, use it.  A reader
%      should be able to find the result in the modern literature under the
%      word used here.
%   2. Where the original term is itself the point -- because it is a coinage
%      Gabriel is introducing, or because it names a notion that later
%      disappeared -- keep an accurate English rendering and give the French
%      original in parentheses at its first occurrence, and thereafter in the
%      terminology table below.
%   3. Sentences are translated, not paraphrased.  Definitions keep their
%      original form (data / axioms as bullet lists) so that the translation
%      can be checked against the original clause by clause.
%
% The table below is the single place where these decisions are recorded.

\begin{table}[htbp]
  \centering
  \small
  \begin{tabular}{p{0.34\textwidth} p{0.28\textwidth} p{0.28\textwidth}}
    \hline
    \textbf{French (original)} & \textbf{English (this translation)} & \textbf{Modern usage}\\[2pt]
    \hline
    univers                     & Grothendieck universe       & Grothendieck universe\\
    ensemble                    & set                          & set\\
    catégorie                   & category                     & category\\
    objet                       & object                       & object\\
    morphisme                   & morphism                     & morphism\\
    morphisme identique         & identity morphism            & identity map\\
    application de composition & composition                  & composition\\
    morphisme functoriel        & natural transformation       & natural transformation\\
    isomorphisme                & isomorphism                 & isomorphism\\
    fondeur / foncteur          & functor                      & functor\\
    Yoncteur                    & Yoneda embedding             & Yoneda embedding\\
    source / but                & source / target              & source / target\\
    mono\-morphisme             & monomorphism                 & monomorphism\\
    épi\-morphisme              & epimorphism                  & epimorphism\\
    sous-objet                  & subobject                    & subobject\\
    quotient                    & quotient                     & quotient\\
    sous-catégorie              & subcategory                 & subcategory\\
    sous-catégorie pleine       & full subcategory             & full subcategory\\
    catégorie duale             & dual category                & dual category\\
    catégorie produit           & product category             & product category\\
    indice                      & index set                    & indexing family\\
    système inductif            & inductive system             & directed system\\
    limite inductive            & inductive limit              & direct limit\\
    système projectif           & projective system            & inverse system\\
    limite projective            & projective limit             & inverse limit\\
    générateur                  & generator                    & generator\\
    générateurs                 & generators                   & generators\\
    ajouté                      & coproduct                    & coproduct\\
    produit                     & product                     & product\\
    \textsc{Ens}                & $\mathbf{Ens}$ (sets)        & $\mathbf{Set}$\\
    \textsc{Ab}                 & $\mathbf{Ab}$ (abelian groups) & $\mathbf{Abel}$\\
    noethérien (à droite)       & (right) noetherian           & noetherian\\
    noethérien                  & noetherian                   & noetherian\\
    localement noethérien       & locally noetherian          & locally noetherian\\
    dimension de Krull          & Krull dimension             & Krull dimension\\
    module pseudo-compact       & pseudocompact module         & pseudocompact module\\
    enveloppe injective         & injective envelope           & injective envelope\\
    foncteur exact à gauche    & left exact functor           & left exact functor\\
    catégorie localisante       & localizing subcategory      & localizing subcategory\\
    catégorie quotient          & quotient category            & quotient category\\
    langage modulo $\mathcal{C}$& language modulo $\mathcal{C}$ & Serre's $\bmod \mathcal{C}$\\
    recollement                 & gluing                       & gluing / recollement\\
    préschéma                   & prescheme                    & scheme / prescheme\\
    faisceau quasi cohérent     & quasi-coherent sheaf        & quasi-coherent sheaf\\
    \hline
  \end{tabular}
  \caption{Translation of the terminology.  Page references are to the
  original French edition, \emph{Bulletin de la Société Mathématique de
  France} \textbf{90} (1962), 323--448.}
  \label{tab:terminology}
\end{table}